\section{Motivation}

    The dynamics of switched affine systems can be described by a state equation of the form 
    \highlight{$\dot{x}(t) = A_{\sigma(t)} x(t) + b_{\sigma(t)}, t \in \mathbb{R}_+$}, where 
    \highlight{$x(t) \in \mathbb{R}^n$} is the state vector and 
    \highlight{$\sigma(t) \in$} \highlight{$\{1, 2, ..., M\}$} is the switching state associated with the matrices \highlight{$A_{\sigma(t)} \in \mathbb{R}^{n \times n}$} 
    and the vectors \highlight{$b_{\sigma(t)} \in \mathbb{R}^{n}$} at each time \highlight{$t$}.

    There is extensive research in the synthesis of controllers for switched affine systems. 
    However, the design procedures are typically concerned with driving the state \highlight{$x(t)$} to 
    the vicinity of a given target, without considering the subsequent evolution of the 
    state in that region.  This issue is especially relevant when the switching state \highlight{$\sigma(t)$} 
    is subjected to dwell-time constraints. In that case, rather than 
    driving the state to an equilibrium point, the design should be aimed at achieving a 
    suitable limit cycle.

    This problem has been addressed, for instance, by \citeonline{benmiloud2018} who uses a 
    hybrid Poincar� map approach to stabilize a switched affine system into the desired 
    limit cycle and a DC-DC converter was employed as an illustrative example. 
    \citeonline{egidio37_discreto} proposed the control design of a state-dependent 
    switching function for discrete-time switched affine systems to stabilize a limit cycle
    determined according to criteria of interest related to the steady-state behavior.
    The control design applies a switched rule to orchestrate the state trajectories based on a time-varying 
    convex Lyapunov function that drives the system from the initial condition toward the desired limit cycle
    The proposed method uses a voltage regulation of a simulated DC-DC multicellular converter as an illustrative example.
    This work was extended to continuous time switched-affine systems by \citeonline{egidio36_continuo}
    

    Some studies form the basis for the present work: 

    \citeonline{marcio2013} solved the problem using predictive control to account for the 
    dwell-time constraints of switching affine systems. 
    Here the controller drove the system to a target state using MILP (Mixed-Integer Linear 
    Programming) and the cyclic trajectory was a consequence of the states as the system 
    approached the target.

    \citeonline{patino2010predictive} defined the target trajectory using nonlinear 
    programming and used this target trajectory in the predictive control in terms of errors.

    \citeonline{matheus2018} proposed a solution with a controller law that drives 
    the switched affine system with dwell-time constraints to a target cyclic trajectory. In 
    this work, the problem is addressed in two stages: first, the problem is interpreted as 
    a convex optimization where the author can find a targeted cyclic trajectory that 
    minimizes a cost function associated with the system performance; then a second part 
    that uses the target trajectory as input and computes perturbation at the switching 
    instants of the reference control signal using predictive control. 




\section{Contributions of the present work}

    The present work proposes an extension of the solution from \citeonline{matheus2020_artigo} to a general case
    of switched affine systems with changing dynamics at switching times and applies the 
    dwell-time constraints using predictive control. 
    
    The work proposed by \citeonline{matheus2020_artigo} has an approach with three basic 
    elements: the determination of a target cyclic trajectory, a linearization of the cycle-to-cycle
    dynamics of the system around the target trajectory that applies the control in the switching time perturbation, 
    the use of the linearized model with predictive formulation to impose the dwell time constraints. This formulation 
    of the solution also has the convenience of solving the problem with quadratic programming convex optimization.

\section{Outline of the remaining chapters}
    
    The remainder of this text is organized as follows. Chapter 2 describes the problem of finding a target periodic trajectory for the 
    given switched affine system. Chapter 3 proposes a predictive control law that drives the system to the target cyclic trajectory. 
    Finally, chapter 4 brings some remarks and considerations of the work.